\chapter{Регистры и счётчики}
\label{ch:regs}

Из прошлой главы мы знаем: один триггер хранит один бит, а тревогу
охранной панели можно запомнить и сбросить ключом. Но панели нужно больше:
хранить код замка, \emph{считать, сколько секунд звучит сирена}, и заставить
её «пищать» раз в секунду, хотя тактовый генератор платы работает на 27~МГц.
Вопрос главы: что получится, если соединить много триггеров, и как из них
сделать хранилище, сдвиг и счёт?

\begin{buildbox}
В Logra откройте Обучение, главу «Счётчики» книги «из блоков», и пройдите
урок «Счётчик одним блоком». Затем урок «Делитель частоты на 2». На вкладке
\term{Диаграммы} посмотрите на все разряды счётчика одновременно: каждый
следующий меняется вдвое реже предыдущего. Запишите, на каком числе счётчик
возвращается в ноль и почему. Ниже мы разберём, как это устроено.
\end{buildbox}

\newcommand{\dff}[3]{% имя, x, y  -- D-триггер 12x14 мм
  \draw[thick, fill=green!8] (#2,#3) rectangle ++(1.2,1.4);
  \node[font=\sffamily\footnotesize] at (#2+0.22,#3+1.05) {D};
  \node[font=\sffamily\footnotesize] at (#2+0.98,#3+1.05) {Q};
  \draw[thick] (#2,#3+0.2) -- (#2+0.18,#3+0.35) -- (#2,#3+0.5);
  \coordinate (#1-d) at (#2,#3+1.05); \coordinate (#1-q) at (#2+1.2,#3+1.05);
  \coordinate (#1-c) at (#2,#3+0.35);}

\section{Регистр хранения}

\term{Регистр} --- это группа D-триггеров с общим тактовым сигналом. По
фронту все триггеры одновременно запоминают свои входы, и регистр хранит
$n$-разрядное число до следующего фронта. На схемах регистр обозначают
буквами RG.

\begin{figure}[H]
\centering
\begin{tikzpicture}
  \foreach \i in {0,...,3} {
    \pgfmathsetmacro{\x}{(3-\i)*3.0}
    \dff{r\i}{\x}{0}
    \draw[wire] (r\i-d) -- ++(-0.4,0) |- ($(r\i-d)+(-0.4,1.0)$) node[above] {$d_\i$};
    \draw[wire] (r\i-q) -- ++(0.5,0) -- ++(0,1.0) node[above] {$q_\i$};
    \draw[wire, FpgaBlue] (r\i-c) -- ++(-0.3,0) -- ++(0,-0.75);
  }
  \draw[wire, FpgaBlue] (-1.4,-0.4) node[left, text=FpgaDark] {CLK} -- (9.0,-0.4);
  \foreach \i in {0,...,3} { \pgfmathsetmacro{\x}{(3-\i)*3.0-0.3} \fill[FpgaBlue] (\x,-0.4) circle (1.3pt); }
\end{tikzpicture}
\caption{Четырёхразрядный регистр: четыре D-триггера с общим тактом}
\label{fig:reg}
\end{figure}

\begin{figure}[H]
\centering
\begin{tikztimingtable}[timing/slope=0.08, xscale=1.5, yscale=1.4, timing/font=\sffamily\small, timing/rowdist=3.2ex]
  CLK      & 8{1L 1H} \\
  $D[3:0]$ & 2D{5} 1.5D{9} 1D{7} 1.5D{3} 4D{C} 6D{0} \\
  $Q[3:0]$ & 1D{?} 2D{5} 2D{9} 2D{3} 4D{C} 5D{0} \\
\end{tikztimingtable}
\caption{Регистр «фотографирует» значение шины $D$ в моменты передних фронтов
CLK. Значение 7 продержалось на входе меньше такта и пришлось между
фронтами --- регистр его просто не заметил}
\label{fig:reg_t}
\end{figure}

Регистры --- это «переменные» аппаратуры. В процессоре в регистрах хранятся
операнды и результаты вычислений; в любом цифровом устройстве --- текущее
состояние, счётчики, данные между этапами обработки. Чтобы регистр менял
значение не на каждом фронте, а только по команде, используют триггеры с
входом разрешения CE (рис.~\ref{fig:ce}).

\begin{diybox}[Регистр хранения]
\textbf{Задание.} \stars{1} Урок «Регистр на четыре бита». \stars{2} Урок
«Регистр с разрешением записи», затем «Регистр со сбросом»: загрузите число,
запретите запись и убедитесь, что оно не меняется. \stars{3} Сохраните в
регистре код замка (четыре бита) и покажите его на четырёх светодиодах.\\
\textbf{В Logra:} Обучение, Конструктор, Диаграммы.\\
\textbf{Ожидаемый результат:} выход обновляется только на фронте, и только
если разрешена запись.\\
\textbf{Как проверить:} на Диаграммах значение 7 на входе, пришедшее между
фронтами, не появилось на выходе (как на рис.~\ref{fig:reg_t}).
\end{diybox}

\subsection{Регистр на Verilog}
\label{sec:regs-reg-v}

На Verilog регистр --- это \code{reg} шириной в несколько бит и один
тактируемый блок. Один и тот же код описывает любое число триггеров, ширину
задаёт параметр \code{W}.

\vfile{\codepath/b8_registers/reg_ce.v}{code/b8\_registers/reg\_ce.v}

\textbf{Как читается.} Заголовок \code{\#(parameter W = 4)} задаёт ширину по
умолчанию, которую можно изменить при подключении модуля. \code{[W-1:0]} ---
диапазон битов от старшего к младшему. По фронту: если \code{rst}, записать
нули (\code{\{W\{1'b0\}\}} --- «повторить \code{1'b0} $W$ раз»), иначе если
\code{we}, записать \code{d}; иначе регистр хранит. Порядок веток --- это
приоритет: сброс сильнее записи.

Тестбенч \code{tb\_reg\_ce.v} проверяет пять тактов: сброс; запись \code{A};
\code{we=0} (регистр хранит \code{A}, хотя на входе \code{5}); запись \code{7};
одновременные \code{rst} и \code{we} (побеждает сброс). \textbf{Как запустить
в Logra:} вкладка \term{Код}, вставьте модуль, Ctrl+Enter; \term{RTL} покажет
четыре триггера с общим тактом и логику выбора перед ними; на \term{Диаграммах}
подайте на \code{d} значения \code{5, A, 7} между фронтами и убедитесь, что
между фронтами выход не меняется.

\begin{diybox}[Регистр на Verilog: допиши модуль]
\textbf{Задание.} \stars{1} Сделайте регистр на 8 бит, не меняя тела модуля:
только параметр. \stars{2} Добавьте второй выход \code{changed}, который равен 1
в такте, когда регистр действительно получил новое значение. \stars{3} Сделайте
регистр, который хранит код замка (4 бита) и только сравнивает: выход
\code{open} равен 1, если на вход подан тот же код.\\
\textbf{Ожидаемый результат:} при $W=8$ на \term{RTL} восемь триггеров.
\end{diybox}

\textbf{Образец ответа.}
\begin{vcode}
module code_lock (
    input        clk,
    input        we,         // запись нового кода
    input  [3:0] code_in,
    output       open        // 1, если code_in совпадает с хранимым кодом
);
    reg [3:0] stored = 4'b0000;
    always @(posedge clk) if (we) stored <= code_in;
    assign open = (code_in == stored);   // сравнение: комбинационная логика
endmodule
\end{vcode}
Первое задание решается строкой \code{reg\_ce \#(.W(8)) r8 (\ldots);} при
подключении модуля в верхнем модуле.

\section{Сдвиговый регистр}

Соединим триггеры цепочкой: выход каждого подаётся на вход следующего. По
каждому фронту содержимое \term{сдвигается} на один разряд, а в первый
триггер записывается новый бит с входа $D_\text{in}$.

\begin{figure}[H]
\centering
\begin{tikzpicture}
  \foreach \i in {0,...,3} {
    \pgfmathsetmacro{\x}{\i*2.3}
    \dff{s\i}{\x}{0}
    \draw[wire, FpgaBlue] (s\i-c) -- ++(-0.3,0) -- ++(0,-0.75);
  }
  \draw[wire] (-1.0,1.05) node[left] {$D_\text{in}$} -- (s0-d);
  \foreach \i [evaluate=\i as \j using int(\i+1)] in {0,1,2} {
    \draw[wire] (s\i-q) -- (s\j-d);
    \fill ($(s\i-q)+(0.3,0)$) circle (1.3pt);
    \draw[wire] ($(s\i-q)+(0.3,0)$) -- ++(0,1.0) node[above] {$q_\i$};
  }
  \draw[wire] (s3-q) -- ++(0.7,0) node[right] {$q_3 = D_\text{out}$};
  \draw[wire, FpgaBlue] (-1.0,-0.4) node[left, text=FpgaDark] {CLK} -- (6.6,-0.4);
\end{tikzpicture}
\caption{Четырёхразрядный сдвиговый регистр}
\label{fig:shift}
\end{figure}

\begin{figure}[H]
\centering
\begin{tikztimingtable}[timing/slope=0.08, xscale=1.5, yscale=1.4, timing/font=\sffamily\small, timing/rowdist=3.2ex]
  CLK            & 8{1L 1H} \\
  $D_\text{in}$  & 2H 2L 12L \\
  $q_0$          & 1L 2H 13L \\
  $q_1$          & 3L 2H 11L \\
  $q_2$          & 5L 2H 9L \\
  $q_3$          & 7L 2H 7L \\
\end{tikztimingtable}
\caption{Единица, поданная на вход, продвигается по регистру на один разряд
за такт}
\end{figure}

\subsection{Сдвиговый регистр на Verilog}
\label{sec:regs-shift-v}

Цепочка триггеров описывается одной строкой: склейка \code{\{\ldots\}}
отбрасывает старший бит и дописывает новый справа.

\vfile{\codepath/b8_registers/shift_sipo.v}{code/b8\_registers/shift\_sipo.v}

\textbf{Как читается.} \code{q[W-2:0]} --- все биты, кроме старшего;
\code{\{q[W-2:0], din\}} дописывает \code{din} в младший разряд. После
присваивания старый старший бит пропадает, остальные «сдвинулись» на разряд
влево. Это тот же регистр, что на рис.~\ref{fig:shift}, только разряды
нумеруются слева направо, а не справа налево.

Обратная задача --- загрузить слово целиком и выдавать по биту:

\vfile{\codepath/b8_registers/shift_piso.v}{code/b8\_registers/shift\_piso.v}

\textbf{Как читается.} Перед каждым триггером стоит мультиплексор «загрузка
или сдвиг»: при \code{load} в регистр попадает слово \code{d}, иначе
содержимое сдвигается влево и справа добавляется 0. Выход \code{dout} ---
старший бит, поэтому слово уходит старшим битом вперёд.

\textbf{Типичная ошибка: «=» вместо «<=».} Посмотрите на два варианта
сдвига, которые отличаются одним знаком:

\vfile{\codepath/b8_registers/shift_bad.v}{code/b8\_registers/shift\_bad.v}

\vfile{\codepath/b8_registers/shift_ok.v}{code/b8\_registers/shift\_ok.v}

В первом блоке присваивания \emph{блокирующие}: строка \code{q1 = q0} уже
видит новое значение \code{q0}, а строка \code{q2 = q1} --- новое \code{q1},
поэтому бит с входа \code{d} попадает во все три выхода за один такт.
Получился не сдвиговый регистр, а один триггер, подключённый к трём выходам.
Во втором блоке присваивания \emph{неблокирующие}: все правые части
вычисляются по значениям, которые были \emph{до} фронта, и триггеры
обновляются одновременно: бит продвигается на один разряд за такт.
Тестбенч \code{tb\_shift.v} подаёт единицу на один такт и проверяет оба
результата: \code{111} у \code{shift\_bad} сразу, а у \code{shift\_ok} ---
\code{001}, затем \code{010} и \code{100}. Он же собирает слово \code{1011}
из потока и выдаёт его обратно из \code{shift\_piso}.

\textbf{Как запустить в Logra.} Вкладка \term{Код}: вставьте
\code{shift\_ok}, Ctrl+Enter, \term{Диаграммы}: подайте \code{d=1} на один
такт и проследите единицу по \code{q0}, \code{q1}, \code{q2}. Затем замените
\code{<=} на \code{=} (или вставьте \code{shift\_bad}) и повторите: на
\term{RTL} вместо цепочки из трёх триггеров будет один триггер с тремя
выходами.

\begin{diybox}[Сдвиг на Verilog: допиши модуль]
\textbf{Задание.} \stars{1} Сделайте сдвиг вправо (новый бит входит в старший
разряд). \stars{2} Добавьте вход разрешения \code{en}: при \code{en=0} регистр
хранит. \stars{3} Сделайте сдвиговый регистр с задержкой на 5 тактов:
выход равен входу, пришедшему 5 тактов назад.\\
\textbf{Ожидаемый результат:} при сдвиге вправо первый пришедший бит
оказывается в младшем разряде, поэтому слово нужно подавать младшим битом
вперёд; задержка на 5 тактов использует пять триггеров.
\end{diybox}

\textbf{Образец ответа.}
\begin{vcode}
module shift_right (input clk, rst, en, din, output reg [3:0] q);
    always @(posedge clk) begin
        if (rst)     q <= 4'b0000;
        else if (en) q <= {din, q[3:1]};   // новый бит слева, младший уходит
    end
endmodule

module delay5 (input clk, din, output dout);
    reg [4:0] sr = 5'b00000;
    always @(posedge clk) sr <= {sr[3:0], din};
    assign dout = sr[4];                    // вход, задержанный на 5 тактов
endmodule
\end{vcode}

\subsection{Зачем сдвигать}

\begin{description}
  \item[Последовательный $\to$ параллельный код.] Биты приходят по одному
    проводу друг за другом (так передают данные по длинным кабелям и
    беспроводным каналам: проводов меньше). После
    $n$ тактов все $n$ битов оказываются в регистре одновременно и могут быть
    прочитаны как число. \lookahead{так работают последовательные
    интерфейсы, глава~\ref{ch:uart}.}
  \item[Параллельный $\to$ последовательный код.] Обратная задача при передаче:
    число загружается в регистр целиком, а затем «выталкивается» по одному
    биту за такт. Для загрузки перед каждым триггером ставят мультиплексор
    «сдвиг/загрузка».
  \item[Задержка сигнала] на заданное число тактов.
  \item[Умножение и деление на 2]: сдвиг двоичного числа влево на один разряд
    умножает его на 2, вправо --- делит.
\end{description}

\begin{figure}[H]
\centering
\begin{tikzpicture}[bit/.style={draw=FpgaBlue, fill=FpgaLight, minimum size=8mm, font=\ttfamily}]
  \node[font=\sffamily\small, anchor=east] at (-0.6,0) {передатчик};
  \foreach \b [count=\i from 0] in {1,0,1,1} \node[bit] at (\i*0.8,0) {\b};
  \draw[very thick, FpgaOrange, -{Stealth}] (3.0,0) -- node[above, font=\sffamily\small, text=FpgaDark] {1 провод: \code{1}, \code{1}, \code{0}, \code{1}\ldots} (7.4,0);
  \foreach \b [count=\i from 0] in {1,0,1,1} \node[bit] at (7.8+\i*0.8,0) {\b};
  \node[font=\sffamily\small, anchor=west] at (10.8,0) {приёмник};
  \node[lbl] at (1.2,-0.8) {параллельный $\to$ послед.};
  \node[lbl] at (9.0,-0.8) {послед. $\to$ параллельный};
\end{tikzpicture}
\caption{Передача числа по одному проводу с помощью двух сдвиговых регистров}
\end{figure}

\subsection{Кольцевой регистр и генератор псевдослучайных чисел}

Если выход последнего триггера соединить со входом первого, получится
\term{кольцевой регистр}: записанная в него единица будет бесконечно бегать по
кругу, как «бегущий огонёк» гирлянды. \lookahead{готовый пример «бегущего
огня» --- в главе~\ref{ch:examples}.}

А если на вход подать «Исключающее ИЛИ» от нескольких разрядов, получится
\term{регистр сдвига с линейной обратной связью} (LFSR). При правильном выборе
отводов он перебирает все $2^n-1$ ненулевых состояний в «перемешанном»
порядке. Такие регистры используются как простые генераторы
псевдослучайных чисел, в помехоустойчивом кодировании и для шифрования.

\begin{figure}[H]
\centering
\begin{minipage}{0.55\textwidth}
\centering
\begin{tikzpicture}
  \foreach \i in {0,...,2} {
    \pgfmathsetmacro{\x}{\i*2.0}
    \dff{l\i}{\x}{0}
  }
  \draw[wire] (l0-q) -- (l1-d); \draw[wire] (l1-q) -- (l2-d);
  \node[gXOR, rotate=180] (x) at (2.0,2.8) {};
  \draw[wire] (l2-q) -- ++(0.4,0) |- (x.input 2);
  \fill ($(l1-q)+(0.3,0)$) circle (1.3pt);
  \draw[wire] ($(l1-q)+(0.3,0)$) |- (x.input 1);
  \draw[wire] (x.output) -| ($(l0-d)+(-0.5,0)$) -- (l0-d);
  \foreach \i in {0,...,2} \node[lbl, below=1mm] at ($(l\i-c)+(0.6,-0.4)$) {$q_\i$};
\end{tikzpicture}
\end{minipage}\hfill
\begin{minipage}{0.4\textwidth}
\centering
\small
\begin{tabular}{c|ccc}
\toprule
такт & $q_0$ & $q_1$ & $q_2$ \\
\midrule
0 & 1 & 0 & 0 \\
1 & 0 & 1 & 0 \\
2 & 1 & 0 & 1 \\
3 & 1 & 1 & 0 \\
4 & 1 & 1 & 1 \\
5 & 0 & 1 & 1 \\
6 & 0 & 0 & 1 \\
7 & 1 & 0 & 0 \\
\bottomrule
\end{tabular}
\end{minipage}
\caption{Трёхразрядный LFSR ($q_0 \leftarrow q_1\oplus q_2$) перебирает все семь
ненулевых состояний и повторяется}
\label{fig:lfsr}
\end{figure}

Состояние «все нули» у такого регистра запретное: из него $q_1\oplus q_2=0$,
и регистр навсегда остаётся нулевым.

\subsection{Кольцо и LFSR на Verilog}
\label{sec:regs-lfsr-v}

Кольцевой регистр --- это сдвиг, в котором выход старшего разряда возвращается
на вход младшего:

\vfile{\codepath/b8_registers/ring4.v}{code/b8\_registers/ring4.v}

\textbf{Как читается.} При сбросе записывается \code{0001} (одна единица); на
каждом такте \code{\{q[2:0], q[3]\}} сдвигает всё влево, а старший бит
дописывается справа. После четырёх тактов единица возвращается на место.
Если вывести \code{q} на четыре светодиода, получится «бегущий огонь».

LFSR --- тот же сдвиг, в котором новый бит вычисляется как исключающее ИЛИ
двух разрядов:

\vfile{\codepath/b8_registers/lfsr3.v}{code/b8\_registers/lfsr3.v}

\textbf{Как читается.} Провод \code{fb} --- обратная связь $q_1\oplus q_2$.
При сбросе записывается \code{001} (только $q_0=1$), а не нуль, поэтому
регистр не застревает: он проходит все семь ненулевых состояний из таблицы на
рис.~\ref{fig:lfsr}. Четырёхразрядный вариант отличается только отводами и
проходит 15 состояний:

\vfile{\codepath/b8_registers/lfsr4.v}{code/b8\_registers/lfsr4.v}

Тестбенч \code{tb\_lfsr\_ring.v} проверяет, что период кольца равен 4, у
\code{lfsr3} --- 7, у \code{lfsr4} --- 15; что все состояния внутри периода
различны и ни разу не встретился нуль. \textbf{Как запустить в Logra:}
вкладка \term{Код}, Ctrl+Enter; на \term{Диаграммах} подайте \code{rst=1} на
один такт, затем \code{rst=0} и выпишите состояния \code{q}: сверьте с таблицей
рис.~\ref{fig:lfsr}. Для четырёх светодиодов выведите \code{q} на выводы платы.

\begin{diybox}[Кольцо и LFSR на Verilog: допиши модуль]
\textbf{Задание.} \stars{1} Сделайте кольцо на 6 разрядов. \stars{2} Сделайте
счётчик Джонсона: на вход младшего разряда подаётся \emph{инверсия} старшего
(\code{\{q[2:0], \textasciitilde q[3]\}}); сколько состояний получилось у 4
разрядов? \stars{3} Что произойдёт с \code{lfsr4}, если убрать сброс и не
задать начальное значение (регистр стартует с \code{0000})?\\
\textbf{Ожидаемый результат:} счётчик Джонсона проходит 8 состояний
(\code{0000, 0001, 0011, 0111, 1111, 1110, 1100, 1000}).
\end{diybox}

\textbf{Образец ответа.}
\begin{vcode}
module johnson4 (input clk, rst, output reg [3:0] q);
    always @(posedge clk) begin
        if (rst) q <= 4'b0000;
        else     q <= {q[2:0], ~q[3]};   // инверсия старшего бита на вход
    end
endmodule
\end{vcode}
У регистра из третьего задания все разряды равны нулю, и обратная связь
$q_3\oplus q_2=0$ навсегда держит его в \code{0000}.

\begin{diybox}[Сдвиг, кольцо и LFSR]
\textbf{Задание.} \stars{1} Урок «Сдвиговый регистр блоком». \stars{2} Уроки
«Последовательный ввод» и «Сборка слова из потока»: проверьте, в каком порядке
(старшим или младшим вперёд) биты попадают в слово. \stars{3} Уроки «Кольцевой
счётчик», «Счётчик Джонсона» и «Псевдослучайные числа: LFSR»: убедитесь, что
четырёхразрядный LFSR проходит 15 состояний.\\
\textbf{В Logra:} Обучение, Конструктор, Диаграммы; для любителей
пример \code{lfsr6}.\\
\textbf{Ожидаемый результат:} единица продвигается на один разряд за такт;
кольцо возвращает её в исходное положение через $n$ тактов.\\
\textbf{Как проверить:} Диаграммы: выпишите состояния и сверьте с таблицей
на рис.~\ref{fig:lfsr}.
\end{diybox}

\section{Счётчики}

\term{Счётчик} --- регистр, значение которого по каждому такту
увеличивается (или уменьшается) на единицу. Счётчики измеряют время, делят
частоту, перебирают адреса памяти.

\subsection{Асинхронный счётчик}

Посмотрите, как меняются разряды при счёте $0, 1, 2, \ldots, 7$: младший бит
переключается каждый шаг, следующий --- в два раза реже, следующий --- ещё в
два раза реже. Каждый разряд переключается, когда предыдущий переходит из 1
в 0. Значит, можно соединить T-триггеры цепочкой: тактовый вход каждого
следующего триггера подключить к \emph{инверсному} выходу предыдущего.

\begin{figure}[H]
\centering
\begin{tikzpicture}
  \foreach \i in {0,...,2} {
    \pgfmathsetmacro{\x}{\i*2.8}
    \draw[thick, fill=green!8] (\x,0) rectangle ++(1.2,1.4);
    \node[font=\sffamily\footnotesize] at (\x+0.22,1.05) {T};
    \node[font=\sffamily\footnotesize] at (\x+0.98,1.05) {Q};
    \node[font=\sffamily\footnotesize] at (\x+0.98,0.35) {$\overline{\text{Q}}$};
    \draw[thick] (\x,0.2) -- (\x+0.18,0.35) -- (\x,0.5);
    \draw[wire] (\x,1.05) -- ++(-0.3,0) node[left, font=\footnotesize] {1};
    \draw[wire] (\x+1.2,1.05) -- ++(0.4,0) -- ++(0,0.8) node[above] {$q_\i$};
  }
  \draw[wire, FpgaBlue] (-1.0,0.35) node[left, text=FpgaDark] {CLK} -- (0,0.35);
  \draw[wire, FpgaRed] (1.2,0.35) -- (2.8,0.35);
  \draw[wire, FpgaRed] (4.0,0.35) -- (5.6,0.35);
\end{tikzpicture}
\caption{Асинхронный (последовательный) трёхразрядный счётчик на T-триггерах}
\label{fig:ripple}
\end{figure}

\begin{figure}[H]
\centering
\begin{tikztimingtable}[timing/slope=0.05, xscale=1.2, yscale=1.4, timing/font=\sffamily\small, timing/rowdist=3.2ex]
  CLK     & 16{1L 1H} \\
  $q_0$   & 1L 7{2H 2L} 2H 1L \\
  $q_1$   & 3L 3{4H 4L} 4H 1L \\
  $q_2$   & 7L 8H 8L 8H 1L \\
  число   & 1D{0} 2D{1} 2D{2} 2D{3} 2D{4} 2D{5} 2D{6} 2D{7} 2D{0} 2D{1} 2D{2} 2D{3} 2D{4} 2D{5} 2D{6} 2D{7} 1D{0} \\
\end{tikztimingtable}
\caption{Каждый следующий разряд меняется вдвое реже: счётчик перебирает числа
$0\ldots7$, а разряд $q_k$ имеет частоту $f_\text{CLK}/2^{k+1}$}
\label{fig:count_t}
\end{figure}

Недостаток асинхронного счётчика --- триггеры переключаются не одновременно,
а «волной», каждый с задержкой относительно предыдущего. В момент перехода
$011 \to 100$ на выходе на короткое время появляются ложные значения
($011\to010\to000\to100$). В длинном счётчике эта «волна» идёт долго.
Поэтому асинхронный счётчик годится лишь там, где его выходы не нужны
остальной схеме одновременно (например, как делитель частоты), и никогда ---
для сравнения значения с константой.

\subsection{Синхронный счётчик}

В синхронном счётчике все триггеры тактируются \emph{одним} сигналом, а новое
значение вычисляет комбинационная схема --- сумматор, прибавляющий единицу.
Все разряды меняются одновременно по фронту. Именно так строятся почти все
счётчики в современных цифровых устройствах.

\begin{figure}[H]
\centering
\begin{tikzpicture}
  \node[draw=FpgaGreen, fill=green!8, thick, minimum width=27mm, minimum height=18mm, font=\sffamily\small, align=center] (reg) at (4.4,0) {регистр\\\footnotesize $n$ бит};
  \node[font=\sffamily\footnotesize, anchor=west] at ([xshift=0.5mm]reg.west) {D};
  \node[font=\sffamily\footnotesize, anchor=east] at ([xshift=-0.5mm]reg.east) {Q};
  \draw[thick] ([xshift=-1.2mm]reg.south) -- ([yshift=1.8mm]reg.south) -- ([xshift=1.2mm]reg.south);
  \draw[wire, FpgaBlue] (reg.south) -- ++(0,-0.8) node[below, text=FpgaDark] {CLK};
  \node[draw=FpgaBlue, fill=FpgaLight, thick, circle, minimum size=12mm, font=\sffamily\bfseries] (add) at (1.4,0) {$+1$};
  \draw[bus, -{Stealth}] (add.east) -- (reg.west);
  \draw[bus, -{Stealth}] (reg.east) -- ++(1.8,0) node[right, text=FpgaDark] {$Q$};
  \draw[bus, -{Stealth}] ($(reg.east)+(0.9,0)$) -- ++(0,1.6) -| (add.north);
  \fill[FpgaBlue] ($(reg.east)+(0.9,0)$) circle (1.8pt);
  \node[lbl] at (4.4,1.9) {обратная связь};
\end{tikzpicture}
\caption{Синхронный счётчик: регистр + сумматор. Каждый такт регистр
записывает $Q+1$}
\label{fig:sync_counter}
\end{figure}

\begin{diybox}[Счётчик]
\textbf{Задание.} \stars{1} Урок «Счётчик одним блоком»: что будет после 15?
\stars{2} Уроки «Счётчик с паузой» и «Счётчик из регистра и «+1»»: сравните
синхронный счётчик, собранный вручную, с готовым блоком. \stars{3} Соберите
счётчик, который считает только при нажатой кнопке.\\
\textbf{В Logra:} Обучение, Конструктор, Диаграммы, RTL; пример
\code{counter}.\\
\textbf{Ожидаемый результат:} после $1111$ счётчик переходит в $0000$.\\
\textbf{Как проверить:} RTL покажет регистр и сумматор «+1», как на
рис.~\ref{fig:sync_counter}.
\end{diybox}

\subsection{Счётчики на Verilog}
\label{sec:regs-cnt-v}

Синхронный счётчик --- это регистр, в который каждый такт записывается
$Q+1$ (рис.~\ref{fig:sync_counter}). На Verilog это одна строка:

\vfile{\codepath/b8_registers/cnt_bin.v}{code/b8\_registers/cnt\_bin.v}

\textbf{Как читается.} По фронту: если сброс --- нули; иначе если разрешён
счёт --- \code{q <= q + 1'b1}. Справа читается \emph{старое} \code{q},
слева записывается новое, сумматор «+1» синтезатор ставит сам. Переполнения
нет в смысле ошибки: сумма отбрасывает перенос, и после \code{1111} получается
\code{0000} (модуль $2^W$). Тестбенч \code{tb\_counters.v} проверяет 16 шагов
подряд, включая переход $15\to0$.

Часто счётчик должен загружать заданное число. Приоритеты задают порядком
\code{if}: сначала сброс, затем загрузка, затем счёт.

\vfile{\codepath/b8_registers/cnt_load.v}{code/b8\_registers/cnt\_load.v}

А реверсивный считает в обе стороны, и «ноль минус один» даёт максимальное
число:

\vfile{\codepath/b8_registers/cnt_updown.v}{code/b8\_registers/cnt\_updown.v}

\textbf{Как читается.} Оператор \code{?:} выбирает между
\code{q + 1} и \code{q - 1}: перед регистром на \term{RTL} вы увидите
сумматор, вычитатель и мультиплексор, а не два отдельных счётчика. Тестбенч
проверяет загрузку (\code{d=9}, потом \code{10}) и переход
$0\to15$ при счёте вниз.

\textbf{Как запустить в Logra.} Вкладка \term{Код}: вставьте
\code{cnt\_bin} и нажмите Ctrl+Enter. На \term{Диаграммах} подайте
\code{rst=1} на один такт, затем \code{en=1} и посмотрите, как каждый следующий
разряд меняется вдвое реже. Повторите для \code{cnt\_load}: подайте
\code{load=1} с \code{d=9} на один такт и убедитесь, что дальше счёт идёт с
девяти.

\begin{diybox}[Счётчик на Verilog: допиши модуль]
\textbf{Задание.} \stars{1} Сделайте счётчик на 8 бит с выходом \code{full},
равным 1 на числе 255. \stars{2} Сделайте счётчик, который считает только
пока \code{en=1} и останавливается (не переходит в 0) на 15. \stars{3}
Сделайте счётчик с шагом \code{step} (3-битный вход): \code{q <= q + step}.\\
\textbf{Ожидаемый результат:} счётчик с остановкой замирает на \code{1111} и
не уходит в 0.
\end{diybox}

\textbf{Образец ответа.}
\begin{vcode}
module cnt_stop (input clk, rst, en, output reg [3:0] q);
    always @(posedge clk) begin
        if (rst)                  q <= 4'd0;
        else if (en && q != 4'd15) q <= q + 4'd1;   // на 15 остаёмся
    end
endmodule

module cnt_step (input clk, rst, input [2:0] step, output reg [7:0] q);
    always @(posedge clk) begin
        if (rst) q <= 8'd0;
        else     q <= q + step;       // ширина step расширяется до 8 бит
    end
endmodule
\end{vcode}

\section{Делитель частоты и счётчик по модулю}

Как видно на рис.~\ref{fig:count_t}, каждый разряд счётчика --- это
делитель частоты на степень двойки. А как поделить на число, не являющееся
степенью двойки, например на 10? Нужен \term{счётчик по модулю $N$}: он считает
$0, 1, \ldots, N-1$ и затем сбрасывается в 0. Для этого к схеме добавляют
компаратор, который сравнивает значение с $N-1$, и мультиплексор, который
вместо $Q+1$ подаёт на регистр 0. Выход компаратора $Q=N-1$ --- это импульс
длиной в один такт, и он появляется ровно раз в $N$ тактов: так получают
сигнал «прошла секунда» из быстрых тактов.

\begin{figure}[H]
\centering
\begin{tikzpicture}
  \node[draw=FpgaGreen, fill=green!8, thick, minimum width=18mm, minimum height=16mm, font=\sffamily\small] (reg) at (5.6,0) {регистр};
  \draw[wire, FpgaBlue] (reg.south) -- ++(0,-0.6) node[right, text=FpgaDark, font=\sffamily\small] {CLK};
  \draw[thick, fill=orange!15] (3.0,0.9) -- (3.8,0.5) -- (3.8,-0.5) -- (3.0,-0.9) -- cycle;
  \node[font=\sffamily\footnotesize] at (3.4,0) {MUX};
  \node[font=\footnotesize, anchor=west] at (3.02,0.5) {0}; \node[font=\footnotesize, anchor=west] at (3.02,-0.5) {1};
  \node[draw=FpgaBlue, fill=FpgaLight, thick, circle, minimum size=10mm, font=\sffamily\bfseries\small] (add) at (1.2,0.5) {$+1$};
  \draw[bus, -{Stealth}] (add.east) -- (3.0,0.5);
  \draw[bus, -{Stealth}] (1.8,-0.5) node[left, text=FpgaDark] {0} -- (3.0,-0.5);
  \draw[bus, -{Stealth}] (3.8,0) -- (reg.west);
  \draw[bus, -{Stealth}] (reg.east) -- ++(1.6,0) node[right, text=FpgaDark] {$Q$};
  \draw[bus] ($(reg.east)+(0.8,0)$) -- ++(0,1.6) -| (add.north);
  \fill[FpgaBlue] ($(reg.east)+(0.8,0)$) circle (1.8pt);
  \node[block, minimum width=20mm] (cmp) at (5.6,-2.9) {$Q = N-1$?};
  \draw[bus, -{Stealth}] ($(reg.east)+(0.8,0)$) |- (cmp.east);
  \draw[arr] (cmp.west) -| node[pos=0.3, below, lbl] {сброс} (3.4,-0.72);
\end{tikzpicture}
\caption{Счётчик по модулю $N$}
\end{figure}

\begin{figure}[H]
\centering
\begin{tikzpicture}
\begin{axis}[
  width=0.82\textwidth, height=5.6cm,
  xlabel={Разрядность счётчика $n$}, ylabel={Частота, Гц},
  ymode=log, ymin=0.1, ymax=3e7, xmin=0, xmax=26, xtick={1,4,8,12,16,20,25},
  label style={font=\sffamily\small}, tick label style={font=\sffamily\small},
  axis lines*=left, grid=major, grid style={FpgaGray!20},
]
  \addplot[very thick, FpgaBlue, mark=*, mark size=1.5pt, samples at={1,2,...,25}] {27e6/2^x};
  \addplot[FpgaRed, thick, dashed, domain=0:26, samples=2] {1};
  \node[font=\sffamily\small, text=FpgaRed, anchor=south west] at (axis cs:1,1.3) {1 Гц --- мигание, заметное глазу};
  \node[font=\sffamily\small, text=FpgaBlue, anchor=south west] at (axis cs:25,0.8) {$n=25$};
\end{axis}
\end{tikzpicture}
\caption{Частота на выходе старшего разряда $n$-битного счётчика при такте
27 МГц: $f = 27\,\text{МГц}/2^n$. Чтобы получить мигание около 1 Гц, нужен
счётчик примерно на 25 бит}
\label{fig:div_plot}
\end{figure}

Но выход старшего бита 25-разрядного счётчика даёт не 1~Гц, а
$27\cdot10^6/2^{25}\approx0{,}8$~Гц. Чтобы получить ровно 1~Гц, нужен счётчик
по модулю $N=27\,000\,000$: он считает от 0 до $26\,999\,999$ (это тоже 25
разрядов) и раз в секунду выдаёт импульс. Если этим импульсом переключать
T-триггер, получится меандр 0,5~Гц (один период за две секунды). Для
меандра ровно 1~Гц берут модуль $13\,500\,000$: импульс приходит каждые полсекунды
и переключает триггер.

\subsection{Счётчик по модулю $N$ и делитель на Verilog}
\label{sec:regs-div-v}

Счётчик по модулю $N$ получается из обычного добавлением сравнения с
$N-1$:

\vfile{\codepath/b8_registers/cnt_mod.v}{code/b8\_registers/cnt\_mod.v}

\textbf{Как читается.} Провод \code{tc} равен 1, пока \code{q} равно $N-1$.
В тактируемом блоке: если сброс \emph{или} \code{tc} --- обнулить, иначе
прибавить единицу. Поэтому значение $N$ не появляется ни в какой момент, а
\code{tc} --- импульс длиной в такт, один раз за $N$ тактов. Параметр
\code{N} можно переопределить при подключении модуля: тестбенч берёт
\code{N = 6} и проверяет, что значения не выходят за 5 и что импульс приходит
ровно 4 раза за 24 такта.

Делитель частоты для платы собирается из таких узлов. Счётчик отсчитывает
$N$ тактов, импульс переключает выход:

\vfile{\codepath/b8_registers/clk_div.v}{code/b8\_registers/clk\_div.v}

\textbf{Как читается.} Для 27~МГц модуль $N=13\,500\,000$ даёт полпериода
0,5~с, то есть меандр 1~Гц на \code{out}; \code{tick} --- импульс «прошло полсекунды».
Счётчик \code{cnt} на 25 разрядов хватает до $2^{25}=33\,554\,432$. Запись
\code{13\_500\_000} (подчёркивания в числе допустимы) удобна для чтения.
Тестбенч \code{tb\_clk\_div.v} подставляет \code{N = 4}: иначе
моделирование потребовало бы 13,5 млн тактов. Он проверяет, что \code{tick}
приходит каждые 4 такта и что период \code{out} равен 8 тактам.

\textbf{Как запустить в Logra.} Вкладка \term{Код}: вставьте \code{cnt\_mod},
Ctrl+Enter, на \term{Диаграммах} подайте сброс на один такт и убедитесь, что
\code{q} проходит $0,\ldots,9$ и возвращается в 0, а \code{tc} поднимается на
девятке. Для \code{clk\_div} задайте параметр \code{N = 4} (в тексте модуля
или в верхнем модуле) --- иначе на Диаграммах 27 миллионов тактов не уместятся
в окно просмотра; настоящее \code{N} оставьте для платы.

\begin{diybox}[Модуль и делитель на Verilog: допиши модуль]
\textbf{Задание.} \stars{1} Сделайте \code{cnt\_mod} с модулем 60 (секунды):
какой минимальной ширины \code{W} хватит? \stars{2} Сделайте делитель, у
которого \code{out} имеет коэффициент заполнения 25\,\% (единица держится четверть
периода). \stars{3} Каскадируйте два \code{cnt\_mod}: единицы секунд (модуль
10) и десятки (модуль 6) с переносом от первого ко второму.\\
\textbf{Ожидаемый результат:} 60 состояний, ширина 6 бит.
\end{diybox}

\textbf{Образец ответа.}
\begin{vcode}
module sec60 (
    input        clk, rst,
    output [3:0] ones,     // единицы секунд 0..9
    output [2:0] tens,     // десятки 0..5
    output       carry     // раз в 60 тактов
);
    reg [3:0] o = 4'd0;
    reg [2:0] t = 3'd0;
    wire o_end = (o == 4'd9);
    wire t_end = (t == 3'd5);
    always @(posedge clk) begin
        if (rst) begin o <= 4'd0; t <= 3'd0; end
        else if (o_end) begin
            o <= 4'd0;
            t <= t_end ? 3'd0 : t + 3'd1;    // десятки меняются, когда единицы закончились
        end else o <= o + 4'd1;
    end
    assign ones  = o;
    assign tens  = t;
    assign carry = o_end & t_end;
endmodule

module duty25 #(parameter N = 4) (input clk, output reg out = 1'b0);
    reg [7:0] cnt = 8'd0;
    always @(posedge clk) begin
        cnt <= (cnt == N - 1) ? 8'd0 : cnt + 8'd1;
        out <= (cnt == 8'd0);               // единица на один такт из N: заполнение 1/N
    end
endmodule
\end{vcode}
Для 60 состояний нужны 6 разрядов, так как $2^5=32<60\le 64=2^6$. Для коэффициента заполнения
25\,\% берут $N=4$.

\begin{diybox}[Модуль $N$ и делитель частоты; этап проекта]
\textbf{Задание.} \stars{1} Урок «Считаем от 0 до 5» (модуль 6) и «Делитель на 4
из счётчика». \stars{2} Уроки «Счётчик со сбросом» и «Счётчик с загрузкой»,
затем «Счётчик на семисегментный индикатор» (глава «Проекты» книги «из
блоков»): сделайте счётчик по модулю 10 с показом на индикаторе.
\stars{3} \textbf{Этап проекта «Охранная панель»: счётчик времени тревоги и
«писк».} Постройте счётчик, который считает секунды, пока тревога из этапа~7
включена, и останавливается на 9. Секунды получите делителем: чтобы
проверка в Диаграммах не длилась 27 миллионов тактов, сначала задайте
малый модуль (например, $N=4$), а настоящий $N=27\,000\,000$ оставьте для платы.
Зуммер «пищит» от бита секунды.\\
\textbf{В Logra:} Обучение (главы «Счётчики», «Проекты»), Конструктор, Диаграммы,
RTL; примеры \code{blink}, \code{countdown}, \code{stopwatch}.\\
\textbf{Ожидаемый результат:} при $N=4$ импульс «секунда» приходит раз в 4
такта; счётчик секунд растёт $0,1,\ldots,9$ и замирает; при сбросе ключом
обнуляется.\\
\textbf{Критерий приёмки:} на Диаграммах измерьте период импульса
(он равен $N$ тактам), выход тревоги на 9-й секунде не возрастает дальше,
а при сбросе все счётчики равны нулю в следующем же такте.
\end{diybox}

\textbf{Этап проекта на Verilog.} Счётчик секунд тревоги получает от делителя
импульс \code{tick} (из \code{clk\_div} или \code{cnt\_mod}) и сигнал
\code{alarm} из главы~\ref{ch:ff}:

\vfile{\codepath/b8_registers/alarm_secs.v}{code/b8\_registers/alarm\_secs.v}

\textbf{Как читается.} Одно условие в одной строке: считать, только если
тревога включена \emph{и} пришёл импульс секунды \emph{и} счёт ещё не дошёл
до 9. Тестбенч \code{tb\_alarm\_secs.v} --- это четыре проверки критерия
приёмки: рост до 3; остановка на 9; ключ обнуляет; при снятой тревоге счёт
не идёт. \textbf{Как запустить в Logra:} вкладка \term{Код}, Ctrl+Enter; на
\term{Диаграммах} подайте \code{alarm=1}, а на \code{tick} --- импульсы
раз в несколько тактов: \code{secs} растёт до 9 и замирает. Чтобы
соединить блоки целиком, сделайте верхний модуль, где \code{alarm\_keep}
(глава~\ref{ch:ff}) даёт \code{alarm}, а \code{clk\_div} с малым \code{N} ---
\code{tick}.

\textbf{Допиши.} Добавьте выход \code{beep}, равный младшему биту
\code{secs}. \emph{Образец:} \code{output beep} и строка
\code{assign beep = secs[0];} --- зуммер «пищит» на нечётных секундах.

\begin{ideabox}
\textbf{Регистр} --- группа триггеров с общим тактом, хранит число.
\textbf{Сдвиговый регистр} перемещает биты по цепочке и преобразует
последовательный код в параллельный и обратно. \textbf{Синхронный счётчик}
--- регистр с сумматором в обратной связи; его разряды делят частоту такта
на степени двойки.
\end{ideabox}

\section{Типичные ошибки}

\begin{warnbox}[Типичные ошибки]
\begin{enumerate}
  \item \textbf{Ошибка на единицу в модуле.} Счётчик по модулю $N$ сбрасывают
        при значении $N-1$, а не $N$. Если сравнивать с $N$, состояний
        окажется $N+1$, а лишнее значение $N$ на мгновение появится на
        выходе.
  \item \textbf{«Бит пролетит весь регистр за один такт».} Все триггеры
        сдвигаются одновременно и берут \emph{старые} значения соседей, поэтому
        бит продвигается ровно на один разряд за фронт.
  \item \textbf{Перепутанный порядок битов.} При приёме потока старший бит
        может идти первым или последним; если слово собирается сдвигом в
        другую сторону, получится число с обратным порядком битов
        ($0011$ превратится в $1100$, а $0001$ в $1000$).
  \item \textbf{Путаница ``мигает'' и ``импульс''.} Старший разряд счётчика
        даёт меандр, а сравнение $Q=N-1$ --- один короткий импульс в такт.
        Для переключения другого узла нужен именно импульс, а не меандр.
  \item \textbf{Нулевое состояние LFSR.} Регистр с исключающим ИЛИ, залитый
        нулями, не выйдет из них никогда; начальное значение должно быть
        ненулевым.
  \item \textbf{Использовать выходы асинхронного счётчика как данные.}
        Из-за «волны» переключения на выходах мелькают ложные числа, и логика,
        прочитавшая их без синхронизации, может сработать лишний раз.
  \item \textbf{Блокирующее присваивание в тактируемом блоке.} Запись
        \code{q1 = q0; q2 = q1;} вместо \code{<=} превращает сдвиговый
        регистр в один триггер: бит «пролетает» все разряды за один такт
        (см. \code{shift\_bad} в разделе~\ref{sec:regs-shift-v}).
  \item \textbf{Ширина счётчика.} \code{reg [3:0] cnt} хранит числа только
        до 15; сравнение \code{cnt == 20} никогда не станет истинным, а
        делитель на 27 миллионов не поместится в 24 разряда ($2^{24}=16\,777\,216$).
  \item \textbf{Нет начального значения.} \code{reg} без сброса и без
        значения при объявлении до первой записи равен \code{X}; счётчик из
        \code{X} так и останется \code{X}.
\end{enumerate}
\end{warnbox}

\section{Вопросы и задачи}
\label{sec:tasks-regs}

Ответы собраны в приложении~\ref{app:answers}.

\begin{taskbox}[Задачи к главе]
\begin{enumerate}
  \item \stars{1}~В четырёхразрядном кольцевом регистре записано $0001$. Какие значения он примет за следующие четыре такта?
  \item \stars{1}~Четырёхразрядный двоичный счётчик находится в состоянии $1111$. Каким станет его состояние после следующего такта? Чему равен модуль счёта?
  \item \stars{2}~Сколько разрядов должен иметь счётчик, чтобы разделить 27~МГц до 1~Гц, если считать до 27\,000\,000?
  \item \stars{1}~До какого значения считает счётчик по модулю 10 и как его сбрасывают?
  \item \stars{2}~Сколько различных ненулевых состояний проходит 4-разрядный LFSR с обратной связью $x^4 + x^3 + 1$ до повторения?
  \item \stars{1}~В чём главное преимущество синхронного счётчика перед асинхронным?
  \item \stars{2}~\textbf{Найди ошибку.} Ученик строил счётчик по модулю 10 и сбрасывал его при значении $1010$. Какую последовательность видит выход? В чём исправить?
  \item \stars{2}~\textbf{Найди ошибку.} Трёхразрядный LFSR из рис.~\ref{fig:lfsr} включили, не задав начальное значение, и он оказался в состоянии $000$. Что будет происходить дальше и как избежать?
  \item \stars{3}~\textbf{Проект.} Счётчик по модулю $N=6$ (или любому другому $N$ по вашему выбору) с показом числа на семисегментном индикаторе. Критерий приёмки: на Диаграммах виден цикл $0,1,\ldots,N-1,0$, значение $N$ не появляется ни в какой момент, а выход «перенос» --- один импульс в такт на числе $N-1$.
  \item \stars{3}~\textbf{Проект (этап 8).} Делитель 27~МГц в 1~Гц: сначала с малым модулем в Диаграммах, затем с настоящим. Критерий приёмки: при модуле $N$ период выходного импульса равен $N$ тактам; для платы модуль $N=27\,000\,000$ и 25 разрядов счётчика.
  \item \stars{3}~\textbf{Проект (этап 8).} Счётчик времени тревоги: считает секунды только пока тревога запомнена (этап~7), останавливается на 9 и обнуляется ключом. Критерий приёмки: четыре проверки на Диаграммах --- рост, остановка на 9, остановка при снятии тревоги, обнуление ключом.
\end{enumerate}
\end{taskbox}

\section{Что дальше}

Теперь у панели есть память, числа и время: код замка в регистре, секунды
тревоги в счётчике и «писк» от делителя. Но всё это пока собрано вручную, и
у схемы нет «режимов»: ключ включил, сигнал пошёл, тревога включилась, код
введён. Как описать такую последовательность состояний строго, чтобы ничего
не потерять? Для этого нужны синхронные схемы и конечные автоматы --- тема
следующей главы.
